\section{Design Twitter}
\label{sec:heap-design-twitter}

\problemstatement{Design a simplified Twitter supporting: $\textit{postTweet}
(\textit{userId}, \textit{tweetId})$ -- a user composes a new tweet;
$\textit{getNewsFeed}(\textit{userId})$ -- return the $10$ most recent
tweet IDs from the given user and everyone they follow, most recent
first; $\textit{follow}(\textit{followerId}, \textit{followeeId})$; and
$\textit{unfollow}(\textit{followerId}, \textit{followeeId})$.}

\begin{patternbox}
\emph{``Most recent $k$ items, drawn from several independently
time-ordered sources''} is exactly Merge K Sorted Lists
(Section~\ref{sec:heap-merge-k-lists}) combined with the Kth Largest
size-bounded heap idea (Section~\ref{sec:heap-kth-largest}), applied
together: each followed user's tweet history is individually
time-ordered (most recent last, or equivalently, sorted by a timestamp),
and we need the top $10$ across all of them -- a direct fusion of two
patterns already established in this chapter, now applied inside a larger
system-design problem.
\end{patternbox}

\subsection*{Understanding the problem carefully}

Each user's tweets, in posting order, already form an individually sorted
sequence (sorted by time posted). Getting a news feed means merging the
relevant followed users' tweet sequences and taking the top $10$ most
recent -- structurally identical to merging $k$ sorted lists and only
needing the first $10$ elements of the merge, rather than the full merge.
We assign a global, strictly increasing \textit{timestamp} to every tweet
at the moment it is posted (a simple counter, incremented on every
$\textit{postTweet}$ call across \emph{all} users), so that ``most
recent'' always corresponds to ``largest timestamp,'' comparable across
different users.

\begin{ideabox}[Data Structures]
\begin{itemize}
  \item A global counter \textit{timestamp}, incremented on every post.
  \item For each user, a list of $(\textit{timestamp}, \textit{tweetId})$
        pairs, in posting order (so the most recent tweet is always at
        the end).
  \item For each user, a set of followee user IDs (a user is always
        considered to follow themselves, for news-feed purposes, without
        needing an explicit \textit{follow} call).
\end{itemize}
\end{ideabox}

\begin{ideabox}[News Feed via a Bounded Max-Heap]
For $\textit{getNewsFeed}(\textit{userId})$: for each followee (including
the user themselves), look only at their \emph{most recent} tweet first
(the tail of their tweet list) -- exactly like the ``current head of each
sorted source'' idea from merging $k$ sorted lists, just read from the end
instead of the front, since within a user's own list, later postings have
larger timestamps. Push each followee's most recent tweet into a max-heap
keyed by timestamp. Repeatedly pop the global maximum, record its tweet
ID, and if that same followee has an earlier tweet remaining (walking
backward through their list), push that earlier tweet in its place. Stop
after popping $10$ tweets or exhausting the heap.
\end{ideabox}

\subsection*{Why this bounded heap pop is correct and efficient}

This is precisely the merge-$k$-sorted-lists technique, with two
adjustments: we traverse each followee's list from most recent to least
recent (so the heap always offers the next-best candidate, by timestamp,
in the right direction), and we stop after $10$ pops rather than draining
the entire merge, since the problem only ever wants the top $10$. Stopping
early is valid precisely because a max-heap's repeated-pop sequence
\emph{is} a fully sorted (by timestamp, descending) traversal of the
union of all followees' tweets -- truncating that sorted traversal at $10$
elements is exactly the $10$ most recent tweets, with no need to ever
materialize the full merge.

\subsection*{Algorithm}

\begin{enumerate}
  \item \textsc{PostTweet}($\textit{userId}, \textit{tweetId}$): append
        $(\textit{timestamp}{+}{+}, \textit{tweetId})$ to
        $\textit{userId}$'s tweet list.
  \item \textsc{Follow}/\textsc{Unfollow}: add/remove
        $\textit{followeeId}$ from $\textit{followerId}$'s followee set.
  \item \textsc{GetNewsFeed}($\textit{userId}$):
  \begin{enumerate}
    \item For each followee $f$ in $\textit{userId}$'s followee set
          (plus $\textit{userId}$ itself): if $f$ has at least one
          tweet, push $(\textit{their last tweet's timestamp}, f,
          \textit{index of that tweet})$ onto a max-heap.
    \item Repeat up to $10$ times: pop the maximum; record its tweet ID;
          if the popped followee has an earlier tweet at
          $\textit{index}-1$, push that one.
    \item Return the recorded tweet IDs, in the order popped (already
          most-recent-first).
  \end{enumerate}
\end{enumerate}

\subsection*{Dry run}

\dryrun{User $1$ posts tweets $5$ (at timestamp $1$) and $3$ (at
timestamp $2$). User $2$ posts tweet $101$ (at timestamp $3$). User $1$
follows User $2$. \textsc{GetNewsFeed}($1$) is called.}

User $1$'s followees (including self): $\{1, 2\}$. Most recent tweet of
user $1$: $(2, 3)$ -- timestamp $2$, tweet $3$. Most recent tweet of user
$2$: $(3, 101)$ -- timestamp $3$, tweet $101$. Heap initialized with both.

\begin{itemize}
  \item Pop max: $(3, 101)$ from user $2$. Record tweet $101$. User $2$
        has no earlier tweet; nothing pushed back.
  \item Pop max: $(2, 3)$ from user $1$. Record tweet $3$. User $1$ has
        an earlier tweet at timestamp $1$ (tweet $5$); push $(1, 5)$.
  \item Pop max: $(1, 5)$ from user $1$. Record tweet $5$. No earlier
        tweet from user $1$.
  \item Heap empty; fewer than $10$ tweets total, so we stop here anyway.
\end{itemize}

The news feed is $[101, 3, 5]$, most recent first -- correct.

\subsection*{C++ implementation}

\begin{lstlisting}
#include <bits/stdc++.h>
using namespace std;

class Twitter {
    int timestamp = 0;
    unordered_map<int, vector<pair<int,int>>> tweets; // userId -> (time, tweetId)
    unordered_map<int, unordered_set<int>> following;  // userId -> followees

public:
    void postTweet(int userId, int tweetId) {
        tweets[userId].push_back({timestamp++, tweetId});
    }

    vector<int> getNewsFeed(int userId) {
        // (timestamp, userId, index into that user's tweet list)
        priority_queue<tuple<int,int,int>> maxHeap;

        auto consider = [&](int uid) {
            if (!tweets[uid].empty()) {
                int idx = tweets[uid].size() - 1;
                maxHeap.push({tweets[uid][idx].first, uid, idx});
            }
        };

        consider(userId);
        for (int f : following[userId]) consider(f);

        vector<int> result;
        while (!maxHeap.empty() && (int)result.size() < 10) {
            auto [time, uid, idx] = maxHeap.top();
            maxHeap.pop();
            result.push_back(tweets[uid][idx].second);

            if (idx > 0) {
                maxHeap.push({tweets[uid][idx-1].first, uid, idx-1});
            }
        }
        return result;
    }

    void follow(int followerId, int followeeId) {
        if (followerId != followeeId) following[followerId].insert(followeeId);
    }

    void unfollow(int followerId, int followeeId) {
        following[followerId].erase(followeeId);
    }
};

int main() {
    Twitter tw;
    tw.postTweet(1, 5);
    tw.postTweet(1, 3);
    tw.postTweet(2, 101);
    tw.follow(1, 2);

    for (int id : tw.getNewsFeed(1)) cout << id << " ";
    cout << endl;
    return 0;
}
\end{lstlisting}

\begin{complexitybox}
\textbf{Time Complexity.} $\textit{postTweet}$, $\textit{follow}$, and
$\textit{unfollow}$ are $O(1)$ amortized. $\textit{getNewsFeed}$ pushes at
most $f$ initial candidates (one per followee, $f \le$ number followed)
and performs at most $10$ pop/push rounds, each $O(\log f)$, giving
\[
  O(f \log f)
\]
per call, independent of the total number of tweets ever posted.
\textbf{Space Complexity.} $O(T)$ to store all $T$ tweets ever posted,
plus $O(f)$ per news feed query for the heap.
\end{complexitybox}

\begin{takeawaybox}
System-design problems frequently turn out to be direct compositions of
patterns you already know: this one is exactly ``merge $k$ sorted
sources'' (Section~\ref{sec:heap-merge-k-lists}) plus ``stop after the
top $k$'' (Section~\ref{sec:heap-kth-largest}), wrapped in bookkeeping for
followers and a global timestamp. Recognizing the composition is more
valuable than treating it as an unfamiliar new problem.
\end{takeawaybox}
